% =============================================================================
% ABSTRACT : English (≈ 350 words, 1/2 page)
% =============================================================================
\chapter*{Abstract}
\addcontentsline{toc}{chapter}{Abstract}

\begin{otherlanguage}{english}
% TODO: finalise after results chapters are stable.

Multi-agent systems built on large language models (LLMs) promise to
orchestrate complex software workflows, yet current hierarchical
architectures (planner--executor, centralised supervisors) suffer from
systemic failure modes documented by the \emph{MAST framework} (Cemri,
2025) and from the diminishing returns observed with frontier models
(Gao, 2025). This thesis investigates an alternative grounded in
\textbf{stigmergy}, the principle of indirect coordination through
modifications of a shared medium theorised by Grass\'e (1959) and
Heylighen (2016) and transposed here to LLM agents. The research
question is: \emph{how can LLM agents coordinate complex tasks by
modifying a shared stigmergic medium rather than through direct
communication or hierarchical control?}

The study follows the \emph{Design Science Research} paradigm
(Hevner, 2004; Peffers, 2007). Five design objectives (DO1 to DO5)
guide the construction of the \textbf{StigmergiAgentic} artifact, a
Python framework articulating a coordination substrate (SQLite-backed
markers with append-only JSONL audit), an ACO-style pressure-driven
decision engine, learning mechanisms (\emph{lessons}, promotion to
\emph{skills}, cross-run protocol persistence), and a governance layer
(\emph{guardrails}, EU~AI~Act Article~14 compliance).

The artifact is empirically validated on two benchmarks. On
\emph{TravelPlanner} (180 constrained planning queries) the C3
stigmergic configuration reaches a final pass rate of 21.1\% with
Gemma and 22.2\% with DeepSeek, significantly outperforming the
\emph{planner-executor} baseline (12.2\%) and marginally outperforming
solo baselines. On \emph{MigrationBench} (Java~8 to Java~17), a
\emph{repair colony} with a closed
\emph{inspect/localize/propose/apply/build/classify/repair} loop
demonstrates a rigorous DSR engineering trajectory. An expert panel
confirms the operational value of the audit trail for regulatory
compliance.

Six generalisable design principles are formulated using the Gregor
(2020) format, foremost among them the \emph{marker primitive as the
sole coordination contract} and \emph{governance through the medium}.
Theoretical contributions extend Malone and Crowston's (1994)
coordination theory to contemporary agentic ecologies.

\vspace{0.3cm}
\noindent\textbf{Keywords}: stigmergy, multi-agent systems, large
language models, Design Science Research, organisational
coordination, code migration, algorithmic governance, EU~AI~Act,
MAST.
\end{otherlanguage}

\cleardoublepage
